% Part: first-order-logic
% Chapter: completeness
% Section: outline

\documentclass[../../../include/open-logic-section]{subfiles}

\begin{document}

\iftag{FOL}
      {\olfileid[fr]{fol}{com}{out}}
      {\olfileid[fr]{pl}{com}{out}}

\olsection{Plan de la démonstration}

La démonstration du théorème de complétude est assez complexe, et il
est facile de s'y perdre à la première lecture. En voici donc le
plan. La première étape consiste à changer de point de vue pour
entrevoir une stratégie. Si l'on formule la complétude ainsi : « si
$\Gamma\Entails !A$, alors $\Gamma\Proves !A$ », il n'est pas
évident de savoir par où commencer. Pour montrer
$\Gamma\Proves !A$, il faut trouver !!a{derivation}, et l'hypothèse
$\Gamma\Entails !A$ ne semble pas nous y aider. Certains systèmes
permettent de construire directement une telle !!{derivation}, mais
nous suivrons une autre voie. Reformulons la complétude : « si
$\Gamma$ est cohérent, alors il est satisfaisable ». Peut-être
pourrons-nous utiliser les informations contenues dans~$\Gamma$ et
sa cohérence pour construire
\iftag{FOL}{!!a{structure}}{!!a{valuation}} qui satisfait chaque
\iftag{FOL}{!!{sentence}}{!!{formula}} de~$\Gamma$. Après tout,
nous savons ce que nous cherchons :
\iftag{FOL}{!!a{structure}}{!!a{valuation}} qui se comporte comme le
décrit~$\Gamma$ !{}

Si $\Gamma$ ne contient que des
\iftag{FOL}{!!{sentence}s atomiques}{!!{propositional variable}s},
il est facile d'en construire un modèle.\iftag{FOL}{ Pour cette
première construction, prenons des !!{sentence}s atomiques sans
égalité, de la forme $\Atom{P}{a_1,\dots,a_n}$, où les $a_i$ sont
des !!{constant}s.}{}
Il suffit de définir
\iftag{FOL}{%
  !!a{domain}~$\Domain{M}$ et une interprétation de~$P$ telle que
  $\Sat{M}{\Atom{P}{a_1,\ldots,a_n}}$. Prenons
  $\Domain{M}=\Nat$ et $\Assign{\Obj c_i}{M}=i$. Pour chaque
  $\Atom{P}{a_1,\ldots,a_n}\in\Gamma$, mettons le tuple
  $\tuple{k_1,\dots,k_n}$ dans $\Assign{P}{M}$, où $k_i$ est
  l'indice de la constante~$a_i$, c'est-à-dire
  $a_i\ident\Obj c_{k_i}$. Dans cette construction initiale,
  $\Assign{P}{M}$ contient exactement les tuples ainsi
  prescrits.}{%
  !!a{valuation}~$\pAssign{v}$ telle que $\pSat{v}{p}$ pour tout
  $p\in\Gamma$. Posons $\pAssign{v}(p)=\True$ si et seulement si
  $p\in\Gamma$.}

Supposons maintenant que $\Gamma$ contienne une !!{formula}
$\lnot !B$, où $!B$ est atomique. On pourrait craindre que la
construction de $\iftag{FOL}{\Struct{M}}{\pAssign{v}}$ nous empêche
de rendre $\lnot !B$ vraie. C'est ici que la cohérence de~$\Gamma$
intervient : si $\lnot !B\in\Gamma$, alors $!B\notin\Gamma$,
sinon $\Gamma$ serait incohérent. Or, si $!B\notin\Gamma$, notre
construction donne
$\iftag{FOL}{\Sat/{M}}{\pSat/{v}}{!B}$, donc
$\iftag{FOL}{\Sat{M}}{\pSat{v}}{\lnot !B}$. Jusque-là, tout va bien.

Que faire si $\Gamma$ contient des formules composées, non atomiques ?
Supposons qu'il contienne $!A\land !B$. Pour rendre cette formule
vraie, nous devons procéder comme si $!A$ et $!B$ appartenaient tous
deux à~$\Gamma$. Si $!A\lor !B\in\Gamma$, il faut rendre au moins
l'une des deux formules vraie, donc procéder comme si l'une d'elles
appartenait à~$\Gamma$.

Voici l'idée : ajoutons des !!{formula}s à~$\Gamma$ de façon
(a) à conserver la cohérence de l'ensemble obtenu, (b) à décider
chaque \iftag{FOL}{!!{sentence}}{!!{formula}}~$!A$ en y plaçant
$!A$ ou $\lnot !A$, et (c) à faire en sorte que, s'il contient
$!A\land !B$, il contienne
aussi $!A$ et $!B$, tandis que, s'il contient $!A\lor !B$, il
contienne au moins l'un des deux, etc. Poursuivons ces ajouts,
éventuellement sans fin, et appelons $\Gamma^*$ l'ensemble formé de
$\Gamma$ et de toutes les !!{formula}s ainsi
ajoutées. La construction précédente fournira alors
\iftag{FOL}{!!a{structure}~$\Struct{M}$}{!!a{valuation}~$\pAssign{v}$}
dont nous pourrons démontrer par récurrence qu'elle satisfait
\iftag{FOL}{tous les !!{sentence}s}{toutes les !!{formula}s}
de~$\Gamma^*$, et donc tout élément de~$\Gamma$ puisque
$\Gamma\subseteq\Gamma^*$. Les conditions (a) et (b) entraînent
les propriétés propositionnelles décrites en~(c).
Un ensemble vérifiant~(b) est dit \emph{complet}. Notre tâche sera
donc d'étendre l'ensemble cohérent~$\Gamma$ en un ensemble cohérent
et complet~$\Gamma^*$.

\iftag{FOL}{%
Une difficulté subsiste. Si $\lexists[x][!A(x)]\in\Gamma$, nous
voudrions choisir une !!{constant}~$c$ et ajouter $!A(c)$. Mais
pourquoi cela serait-il toujours possible ? Notre langage ne contient
peut-être que quelques !!{constant}s, et $\Gamma$ pourrait contenir
$\lnot !A(c)$ pour chacune d'elles. Ajouter aussi $!A(c)$ rendrait
l'ensemble incohérent : nous ne saurions plus si $\Struct{M}$ doit
rendre $!A(c)$ ou $\lnot !A(c)$ vraie. Il se peut même que le langage
de~$\Gamma$ ne possède aucune constante, comme celui de la théorie
des ensembles.

La solution consiste à ajouter d'emblée une infinité de constantes,
ainsi que des !!{sentence}s qui les relient convenablement aux
quantificateurs. Il faudra évidemment vérifier que cela n'introduit
pas d'incohérence.

La construction initiale fonctionne bien lorsque les !!{sentence}s
atomiques ne contiennent que des constantes. Mais le langage peut
aussi comporter des !!{function}s. Il pourrait alors être difficile
de trouver les fonctions sur~$\Nat$ à attribuer à ces symboles pour
que tout fonctionne. Voici une autre idée : au lieu d'interpréter
$\Obj c_i$ par~$i$, prenons comme domaine l'ensemble des constantes
elles-mêmes. La !!{structure}~$\Struct M$ pourra alors attribuer
chaque constante à elle-même : $\Assign{\Obj c_i}{M}=\Obj c_i$.
Prenons pour $\Domain{M}$ l'ensemble de tous les
\emph{termes clos} du langage. Il existe alors une interprétation
naturelle des symboles de fonction par des fonctions qui prennent des
termes en arguments et produisent des termes en valeurs. Au symbole
$\Obj f^n_i$, nous attribuons la fonction qui envoie les $n$ termes
$t_1,\dots,t_n$ sur le terme $\Obj f^n_i(t_1,\dots,t_n)$.

Reste à traiter~$\eq$. Ce !!{predicate} a une interprétation fixée :
$\Sat{M}{\eq[t][t']}$ si et seulement si
$\Value{t}{M}=\Value{t'}{M}$. Si la valeur de chaque terme~$t$ est
le terme~$t$ lui-même, la !!{structure} ne rendra vrai aucun
!!{sentence} de la forme $\eq[t][t']$ à moins que $t$ et $t'$ ne
soient un seul et même terme. Cela pose problème : presque toute
théorie intéressante dans un langage comportant des !!{function}s a
des théorèmes de la forme $\eq[t][t']$ pour des termes distincts ;
par exemple, en arithmétique, $\eq[(\Obj 0+\Obj 0)][\Obj 0]$.
Nous modifions donc le domaine de~$\Struct M$ : ses objets seront
des classes de termes plutôt que des termes isolés. Deux termes
appartiendront à la même classe lorsque les égalités de~$\Gamma^*$
exigent qu'ils aient la même valeur. Si $\Gamma$ est une théorie
arithmétique, l'une de ces classes contiendra par exemple
$\Obj 0$, $(\Obj 0+\Obj 0)$, $(\Obj 0\times\Obj 0)$, etc. Elle
sera l'interprétation de~$\Obj 0$ et la valeur de chacun de ses
termes, notamment $(\Obj 0+\Obj 0)$. Ainsi, l'!!{sentence}
$\eq[(\Obj 0+\Obj 0)][\Obj 0]$ sera vrai dans cette nouvelle
!!{structure}.}{}

Voici les étapes de la démonstration. Nous étudierons d'abord les
ensembles cohérents et !!{complete}s. En particulier, nous montrerons
qu'un tel ensemble contient $!A\land !B$ si et seulement s'il
contient à la fois $!A$ et $!B$, qu'il contient $!A\lor !B$ si et
seulement s'il contient au moins l'une des deux, etc.
(\olref[ccs]{prop:ccs}).\iftag{FOL}{ Nous définirons ensuite les
ensembles saturés. Un tel ensemble contient des implications
qui relient les !!{sentence}s quantifiés à leurs instances
(\olref[hen]{defn:henkin-exp}). Nous montrerons que tout ensemble
cohérent~$\Gamma$ s'étend en un ensemble cohérent et saturé~$\Gamma'$
(\olref[hen]{lem:henkin}). S'il est cohérent, saturé et
!!{complete}, un ensemble vérifie en outre les propriétés suivantes :
il contient
\iftag{prvEx}{$\lexists[x][!A(x)]$ si et seulement s'il contient
$!A(t)$ pour un terme clos~$t$}{}\iftag{defEx,defAll}{}{ et
}\iftag{prvAll}{$\lforall[x][!A(x)]$ si et seulement s'il contient
$!A(t)$ pour tout terme clos~$t$}{}.
(\olref[hen]{prop:saturated-instances}).}{}
Nous partirons alors de l'ensemble cohérent
$\iftag{FOL}{\Gamma'}{\Gamma}$,\iftag{FOL}{ déjà saturé,}{} pour
l'étendre en un ensemble~$\Gamma^*$
\iftag{FOL}{saturé, cohérent et}{cohérent et} !!{complete}
(\olref[lin]{lem:lindenbaum}). C'est cet ensemble qui servira à
définir notre
\iftag{FOL}{modèle des termes~$\Struct M(\Gamma^*)$}{!!{valuation}~$\pAssign v(\Gamma^*)$}.
\iftag{FOL}{Le domaine du modèle des termes est l'ensemble des termes
clos, et l'interprétation de ses !!{predicate}s est déterminée par
les !!{sentence}s atomiques}{La valuation est déterminée par les
!!{propositional variable}s} de~$\Gamma^*$
(\olref[mod]{defn:termmodel}). Nous utiliserons les propriétés des
ensembles cohérents, !!{complete}s\iftag{FOL}{ et saturés}{} pour
montrer que
$\iftag{FOL}{\Sat{M(\Gamma^*)}}{\pSat{v(\Gamma^*)}}{!A}$ si et
seulement si $!A\in\Gamma^*$ (\olref[mod]{lem:truth}), et donc, en
particulier,
$\iftag{FOL}{\Sat{M(\Gamma^*)}}{\pSat{v(\Gamma^*)}}{\Gamma}$.
\iftag{FOL}{Enfin, dans un langage avec~$\eq$, nous définirons
un modèle quotient des termes
(\olref[ide]{defn:term-model-factor}) et montrerons qu'il
satisfait~$\Gamma^*$ (\olref[ide]{lem:truth}).}{}

\end{document}
